\begin{musica}{Clavelitos}{Popular Espanhola}{Tuna de Engenharia da Universidade do Porto}
\nota{Introdução: solo#1:}
\begin{tablatura}
E-|-------------------------------------------|-------------------------------------------------
B-|-------------------------------------------|-------------------------------------------------
G-|-------------------------------------------|-------------------------------------------------
D-|-------------------------------------------|-------------------------------------------------
A-|---------2----------2--0--3--2--0----------|-------------------------------------------------
E-|--0--4------0--4--------------------3--1---|-------------------------------------------------
\end{tablatura}

Mo\ac{Am}{ci}ta \ac{E}{dá} me el cla\ac{Am}{vel},
Da me el clavel de tu \ac{E}{bo}ca.
Que para eso no hay que tener
Mucha verguenza ní \ac{Am}{po}ca.
\ac{G}{Yo} te daré el casca\ac{C}{vél},
\ac{G}{Te} lo prometo, mo\ac{C}{ci}ta;
\ac{E}{Si} tu me dás esa \ac{Am}{mi}el
\ac{E}{Que} llevas en lá bo\ac{A}{qui}ta.

\nota{Refrão:}
Clave\ac{A}{li}tos, clavelitos,
Clavelitos de mí cora\ac{E}{zón};
Hoy te traigo, clavelitos
Colorados igual que un fre\ac{A}{són}.
Si algun dia, clavelitos
No lo\ac{A7}{gra}ra poder te tra\ac{D}{er},
No te \ac{D}{creas} que yá no te \ac{A}{qui}ero,
Es que \ac{E}{no} te los pude co\ac{A E}{ger}. \ac{A}{}_

Lá tarde que a média luz
Vi tu boquita de guinda,
Yo no he visto en Santa Cruz
Otra boquita tan linda.
Y luego al ver el clavel
Que llevavas en el pelo,
Mirando crei yo ver
Un pedacito de cielo.

\nota{Refrão}

\posicoes{C: x32010, D: xx0232, E: 022100, G: 320033, A: x02220, A7: x02020,
  Am: x02210}
\end{musica}
